%==============================================================================
%  CHAPITRE 2 — ÉTAT DE L'ART
%==============================================================================
\titreduchapitre{Synthèse de l'état de l'art}{chap:etat_art}

%==============================================================================
\section{Introduction}
\label{sec:ea-intro}

Le chapitre~\ref{chap:fondements} a fixé le vocabulaire. Il reste à
examiner ce que la littérature établit déjà, et ce qu'elle laisse ouvert,
lorsqu'il s'agit de réguler un apprenant isolé devant une vidéo de cours.

Deux familles de travaux se croisent sans se rejoindre pleinement. D'un
côté, on détecte l'engagement : parfois dans les messages de forum
(présence cognitive, émotions épistémiques), parfois sur le visage, la voix
ou la posture. De l'autre, on intervient : tuteurs intelligents,
règles d'adaptation, puis, depuis peu, texte généré par un grand modèle de
langue. Peu de prototypes enchaînent les deux : estimer un état, choisir une
stratégie, produire un contenu. C'est cet enchaînement que vise
\emolearn{}.

Le chapitre s'organise en six sections : MOOC et engagement apprenant,
détection automatique, psychologie des émotions, IA générative, systèmes
adaptatifs, puis synthèse critique et positionnement. Ce parcours prépare
les choix du chapitre~\ref{chap:conception}.

%==============================================================================
\section{MOOCs et engagement apprenant}
\label{sec:ea-mooc}

Trois questions structurent cette section : les défis de rétention, la
mesure de l'engagement en e-learning, puis les indicateurs émotionnels
et cognitifs.

\subsection{Défis de rétention et d'engagement}
\label{sec:ea-mooc-retention}

\textcite{daniel2012making} et \textcite{breslow2013studying}
ont d'abord décrit le MOOC comme un phénomène
d'accès ouvert à grande échelle. Les chiffres de rétention
sont connus depuis longtemps : une grande partie des inscrits part dans
les premières semaines, et le taux moyen d'achèvement stagne près de
10\,\% \cite{onah2014dropout,jordan2015massive,reich2019mooc}. Onah,
Jordan et Reich établissent ce constat. Kizilcec et Hew montrent que le
manque de temps n'explique pas à lui seul les départs : la baisse
d'engagement y contribue également
\cite{kizilcec2013deconstructing,hew2016promoting}. Le défi de rétention
porte donc aussi sur ce que les inscrits vivent pendant le cours, au-delà
du calendrier et du volume de travail déclaré.

\subsection{Mesure de l'engagement en e-learning}

L'engagement désigne la mobilisation cognitive, comportementale et
émotionnelle de l'apprenant sur une tâche \cite{fredricks2004school}. La
définition aujourd'hui la plus partagée distingue trois composantes :
\begin{itemize}
  \item l'engagement comportemental (participation, persévérance,
  effort observable)~,
  \item l'engagement cognitif (profondeur du traitement, stratégies
  d'apprentissage, présence cognitive)~,
  \item l'engagement émotionnel (intérêt, ennui, confusion,
  frustration, sentiment d'affiliation).
\end{itemize}
En e-learning, on mesure surtout des proxies : connexion, clics,
durée de session, progression, notes. Ils sont utiles, mais incomplets :
un étudiant peut rester connecté et être déjà ailleurs. Les journaux de
plateforme voient le comportement. Les forums donnent un aperçu de la
cognition, à condition que l'étudiant écrive. Compter les clics ne
suffit donc plus : il faut une mesure plus proche de ce que l'apprenant
est en train de vivre
\cite{kizilcec2013deconstructing,dmello2015review}.

\subsection{Indicateurs émotionnels et cognitifs}

L'émotion, pourtant liée à la persévérance
\cite{pekrun2006control,dmello2012dynamics}, n'apparaît pas dans les logs.
\emolearn{} s'appuie sur ces indicateurs, dans l'hypothèse qu'une
régulation plus précoce agit sur la suite du parcours. L'étude de Liu \emph{et al.} sur le forum d'\emph{Introduction
to Psychology} (8\,867 inscrits) va dans ce sens : émotion et cognition
bougent ensemble, et prédisent mieux la réussite lorsqu'on les lit
conjointement \cite{liu2022automated}. Un seul score d'\og engagement \fg{}
efface ces nuances.

Dès que l'on dispose de texte, la revue de Sakeef \emph{et al.} montre que
trois grilles reviennent pour coder l'engagement cognitif des
forums \cite{sakeef2025detecting}. La plus citée est le modèle de
Community of Inquiry \cite{garrison1999critical} : présence sociale,
enseignante et cognitive, cette dernière en quatre phases (déclenchement,
exploration, intégration, résolution). Hu \emph{et al.} et Liu \emph{et al.}
s'en servent pour entraîner des classifieurs
\cite{hu2021automatic,liu2023moocbert}. ICAP, proposé par Chi et Wylie,
classe plutôt le mode d'activité : passif, actif, constructif,
interactif \cite{chi2014icap}. Plus on construit ou on interagit, plus le
traitement est censé être profond.

Ces indicateurs cognitifs aident à relire un fil de discussion
après coup. Ils disent peu d'un visionnage sans contribution
écrite. Dans ce cas, les indicateurs émotionnels visibles, taxonomie
\textsc{DAiSEE} : engagement, ennui, confusion, frustration sont plus
directs : un état, un besoin pédagogique \cite{gupta2016daisee}.

%==============================================================================
\section{Détection automatique de l'engagement}
\label{sec:ea-detection}

On ne peut pas demander à un enseignant de coder des milliers de messages,
ni d'observer chaque visage. D'où l'intérêt d'un estimateur automatique.
Cette section examine trois axes : les approches multimodales (visuelle,
comportementale, physiologique), les corpus publics (\textsc{DAiSEE},
\textsc{EmotiW}, et les ressources complémentaires), puis les performances
rapportées dans la littérature. Les scores d'\emolearn{} lui-même (FER+AU,
WavLM, fusion) sont ceux du chapitre~\ref{chap:validation}, ils ne
constituent pas un palmarès de l'état de l'art.

\subsection{Approches multimodales}
\label{sec:ea-detection-multimodal}

D'Mello, dans sa méta-analyse, insiste sur un point pratique : ennui,
confusion, frustration et engagement reviennent souvent, et mélanger les
modalités aide, à condition de calibrer \cite{dmello2015review}. Trois
familles de signaux structurent cet axe.

\paragraph{Voie visuelle:}
En l'absence de production écrite, le signal disponible est celui que
l'on peut filmer. Whitehill \emph{et al.} ont montré, dès 2014, qu'un
classifieur de visage
peut approcher le jugement humain sur l'engagement
\cite{whitehill2014faces}. Depuis, les réseaux profonds et les Transformers
ont relevé le niveau \cite{li2020deepfer,dosovitskiy2021vit}. En classe,
d'autres travaux suivent le regard ou le décrochage attentionnel
\cite{goldberg2021attentive,hutt2021automated}. Sur \textsc{DAiSEE},
plusieurs architectures hybrides spatio-temporelles ont été testées
\cite{abedi2021improving,gupta2023multimodal}.

Xiong, Guo et Xu proposent un CNN-Transformer qui combine l'information
fine des expressions faciales (CNN) et l'information plus grossière de la
posture (Transformer), en réduisant l'encombrement visuel du fond
\cite{xiong2024cnntransformer}. Sur un jeu de données de classe réelle,
annoté par crowdsourcing puis validation experte, le modèle atteint
92,9\,\% d'accuracy et surpasse plusieurs architectures de deep learning.
Ce travail est le plus proche d'\emolearn{} du point de vue du
signal (vision en situation d'apprentissage). Il s'en distingue
toutefois sur deux points décisifs : il ne fusionne pas la voix, et il ne
relie pas la détection à une génération de contenu. L'identification
s'arrête au diagnostic.

Côté voix: souvent traitée à part, mais bien une modalité perceptive,
le passage des descripteurs acoustiques classiques aux représentations
auto-supervisées (WavLM), évaluées notamment via le
benchmark SUPERB, a transformé la reconnaissance
d'émotions vocales
\cite{chen2022wavlm,pepino2021emotion,yang2021superb,latif2023survey}.
Le fine-tuning de ces encodeurs \cite{chen2023exploring} et les
architectures Transformer réduisent notamment l'écart de performance
sur la valence, longtemps le point faible du SER
\cite{wagner2023dawn}. Ces avancées justifient le choix actuel
d'\emolearn{} : un encodeur WavLM fine-tuné vers les quatre
états \textsc{DAiSEE}, complété par une ancre émotionnelle
\textsc{IEMOCAP}, plutôt que des seuils acoustiques ou un simple mapping
émotion\,$\rightarrow$\,engagement. Côté visage, le runtime retient un
Random Forest FER+AU (sept scores d'un ViT pré-entraîné, plus
Action Units) comme juge produit, un ViT fine-tuné vers les quatre
états filtre les images peu confiantes et sert de mode dégradé, sans
choisir l'état transmis à la fusion.

\paragraph{Voie comportementale:}
On rassemble ici les traces d'activité: clics, durée de session,
progression dans un MOOC, souris, posture, et le texte laissé sur un forum.
Kizilcec et al. montrent que ces proxies séparent des
sous-populations d'inscrits, sans dire ce que l'étudiant est en train de
vivre \cite{kizilcec2013deconstructing}. La posture, chez Xiong, est déjà
un signal comportemental filmé \cite{xiong2024cnntransformer}. Dans les
MOOC, le matériau comportemental le plus riche reste le texte des
forums.

Sakeef \emph{et al.} proposent une revue des cadres (CoI, ICAP, KCSA) et des
méthodes de TALN / apprentissage profond utilisées pour détecter l'engagement
cognitif dans ces messages \cite{sakeef2025detecting}. Ils soulignent que
l'annotation manuelle est trop coûteuse pour un suivi en continu, et que
l'automatisation devient indispensable. Ils rappellent aussi une limite
structurelle : l'absence d'indices visuels et auditifs dans le texte rend
l'inférence de l'état interne partielle. Cette limite motive, pour
\emolearn{}, le recours aux modalités visage et voix (corpus annotés)
plutôt qu'au seul forum.

Hu, Ferreira Mello et Gašević comparent un CNN à des Random Forests
fondés sur des traits linguistiques pour classer la présence cognitive
selon le CoI \cite{hu2021automatic}. Le CNN atteint un $\kappa$ de Cohen de
0,528, comparable aux approches classiques, avec une généralisabilité
inter-cours similaire. L'apport principal est l'explicabilité (XAI) au
niveau des mots. Le travail reste une tâche de catégorisation de messages,
sans lien avec une intervention pédagogique automatique.

Liu \emph{et al.} (2022) introduisent BERT-CNN pour détecter conjointement
l'engagement émotionnel et cognitif dans le forum du MOOC
\emph{Introduction to Psychology}
\cite{liu2022automated}. Par rapport aux modèles antérieurs, le F1
s'améliore d'environ 10\,\% (émotion) et 8\,\% (cognition).
Méthodologiquement, le papier établit qu'un détecteur robuste est un
prérequis à toute régulation, opérationnellement, il s'arrête à la
prédiction de la réussite, sans générer de contenu adaptatif.

MOOC-BERT pousse plus loin le pré-entraînement : le modèle est pré-entraîné
sur le MOOC-Corpus (icourse163.org : 800 cours, cinq disciplines,
3{,}76~Go), puis fine-tuné pour identifier la présence cognitive
\cite{liu2023moocbert}. Il surpasse les algorithmes représentatifs, y
compris en généralisation inter-cours. La boucle s'arrête à l'aide à
l'enseignant : fournir un diagnostic, non produire l'intervention.

Liu \emph{et al.} (2021) articulent analyse de réseaux sociaux et analyse de
réseaux épistémiques sur 3\,925 participants d'un forum MOOC
\cite{liu2021exploring}. Les groupes à haute et basse réussite diffèrent
en centralité, les connexions cognitives sont plus fortes chez les
apprenants performants. Ce résultat concerne une dynamique
collective (forum), distincte du scénario individuel visé par
\emolearn{}.

\paragraph{Voie physiologique:}
Une troisième famille EEG, réponse électrodermale, rythme cardiaque,
parfois ECG existe surtout en laboratoire
\cite{picard1997affective,dmello2015review,abedi2021improving}. Elle sort du
périmètre d'\emolearn{} : le scénario visé est une vidéo de cours, avec
webcam et microphone, sans capteur médical.

\subsection{Datasets publics d'engagement}
\label{sec:ea-detection-datasets}

Les corpus visuels publics d'engagement restent rares. On en
retient ici ce qui structure la détection.

\textsc{DAiSEE} fournit des vidéos d'apprenants annotées en quatre états
(engagement, ennui, confusion, frustration), chacun sur une intensité de
0 à 3, en conditions naturelles (\og dans la nature \fg{})
\cite{gupta2016daisee}. La classe Engagement y domine largement, ce qui
rend l'exactitude globale insuffisante comme seule métrique. C'est la taxonomie retenue par
\emolearn{}, parce que chaque état appelle une réponse pédagogique
distincte.

Le défi \textsc{EmotiW} propose aussi des pistes \og dans la nature \fg{}, dont une
tâche d'engagement étudiant \cite{dhall2018emotiw}. Les étiquettes y sont
surtout des niveaux d'engagement, pas les quatre états
pédagogiques de \textsc{DAiSEE}. Les scores des éditions successives ne se
lisent donc pas comme un classement sur le même problème.

\textsc{IEMOCAP} offre un signal vocal riche et, dans sa version
audiovisuelle, un cadre où visage et voix proviennent du même locuteur
\cite{busso2008iemocap}. Les étiquettes y décrivent l'émotion dyadique,
et non l'engagement d'un apprenant en MOOC. \emolearn{} s'appuie
explicitement sur ces deux ressources complémentaires : \textsc{DAiSEE}
pour la taxonomie d'engagement facial, \textsc{IEMOCAP} pour l'émotion
vocale.

D'autres jeux restent propriétaires : Whitehill \emph{et al.}
annotent leurs propres vidéos d'étudiants \cite{whitehill2014faces}, 
Xiong \emph{et al.} collectent des images de classe faute de dataset
ouvert alliant visage et posture \cite{xiong2024cnntransformer}. Les
travaux sur forums (icourse163.org, bancs CoI) et les essais PAGE ou Kong
n'offrent pas un dataset public d'engagement visuel.

\subsection{État de l'art des performances de détection}
\label{sec:ea-detection-perfs}

Comparer un $\kappa$ sur du texte CoI, une accuracy sur quatre niveaux  \textsc{DAiSEE} et 92,9\,\% sur un corpus de classe privé
n'a de sens que si l'on rappelle que la tâche n'est pas la même.
Le tableau~\ref{tab:ea-perfs-detection} présente les chiffres
effectivement rapportés par les articles déjà discutés. Il ne prétend pas
couronner un vainqueur unique.

{\scriptsize
\setlength{\tabcolsep}{3pt}
\renewcommand{\arraystretch}{1.08}
\begin{xltabular}{\textwidth}{|>{\raggedright\arraybackslash}X|>{\raggedright\arraybackslash}X|>{\raggedright\arraybackslash}X|>{\raggedright\arraybackslash}X|}
\hline
\rowcolor{gray!20}
\textbf{Référence} & \textbf{Tâche / signal} & \textbf{Corpus} &
Résultat clé \\
\hline
\endfirsthead
\hline
\rowcolor{gray!20}
\textbf{Référence} & \textbf{Tâche / signal} & \textbf{Corpus} &
Résultat clé \\
\hline
\endhead
\endfoot
\endlastfoot
\textcite{whitehill2014faces}
  & Visage, 4 niveaux d'engagement
  & Corpus propre (non public)
  & Accord proche des annotateurs humains, \textcite{gupta2016daisee} rappellent $\approx 69\,\%$ pour SVM + Gabor \\
\hline
\textcite{gupta2016daisee}
  & Vidéo, 4 niveaux d'engagement
  & \textsc{DAiSEE} (public)
  & Baselines des auteurs : InceptionNet $46{,}4\,\%$, C3D $56{,}1\,\%$, LRCN $57{,}9\,\%$ \\
\hline
\textcite{abedi2021improving}
  & Vidéo spatio-temporelle, 4 niveaux
  & \textsc{DAiSEE}
  & ResNet+TCN : $63{,}9\,\%$ d'accuracy (mieux que les baselines \textsc{DAiSEE} et plusieurs hybrides) \\
\hline
\textcite{xiong2024cnntransformer}
  & Visage + posture, 4 niveaux
  & Images de classe (propre)
  & Accuracy $92{,}9\,\%$, pas de voix, corpus non public \\
\hline
\textcite{hu2021automatic}
  & Texte, phases CoI
  & Deux cours annotés
  & $\kappa$ de Cohen $= 0{,}528$ (CNN $\approx$ Random Forests) \\
\hline
\textcite{liu2022automated}
  & Forum, émotion + cognition
  & icourse163.org ($n=8\,867$)
  & F1 $+10\,\%$ (émotion) et $+8\,\%$ (cognition) vs modèles antérieurs \\
\hline
\textsc{EmotiW} \cite{dhall2018emotiw}
  & Engagement \og dans la nature \fg{}
  & Pistes du défi (public)
  & Scores d'édition en édition, niveaux d'engagement, non comparables à la taxonomie \textsc{DAiSEE} \\
\hline
\end{xltabular}
\xlcaption{Performances de détection rapportées.}
\label{tab:ea-perfs-detection}
}

Deux lectures s'imposent. Premièrement, sur \textsc{DAiSEE}, même les
meilleures architectures citées ici restent dans une fourchette d'accuracy
modérée ($57{,}9\,\%$ à $63{,}9\,\%$ pour quatre niveaux d'engagement),
parce que la classe haute domine : un classifieur naïf « toujours
engagé » paraît déjà fort. C'est pourquoi \emolearn{} retient, au
chapitre~\ref{chap:validation}, le F1 macro et le rappel des états
urgents plutôt que l'accuracy seule. Deuxièmement, un score élevé sur un
corpus privé (Xiong) ou un $\kappa$ sur du CoI ne se substitue pas à une
évaluation visage--voix vers les quatre états \textsc{DAiSEE}.

Identifier un état ne suffit pas à justifier son importance pédagogique,
ni à en déduire une intervention. La section suivante rappelle les
fondements psychologiques, les deux suivantes traitent de la génération
et des systèmes adaptatifs.

%==============================================================================
\section{Psychologie de l'apprentissage et émotions}
\label{sec:ea-psycho}

Les cadres sur les MOOC et l'engagement indiquent quoi mesurer.
Ils disent moins \emph{pourquoi} un état change la suite du parcours. Trois
points sont examinés : l'impact des émotions, les théories de la charge
cognitive, puis les états d'engagement (ennui, confusion,
frustration, flow).

\subsection{Impact des émotions sur l'apprentissage}

Picard et Pekrun ont, chacun de leur côté, rappelé que l'émotion n'est pas
un bruit autour de la cognition : elle oriente l'attention et la
persévérance \cite{picard1997affective,pekrun2006control}. Les revues
récentes sur les environnements technologiques aboutissent au même constat
\cite{loderer2020beyond,graesser2020emotions}. Une émotion n'est donc pas
un à-côté du cours : elle décide, en partie, si l'étudiant reste ou part.

\subsection{Théories de la charge cognitive}

Sweller décrit le couple sous l'angle de la charge : trop de matière à la
fois, et l'étudiant se perd (surcharge, confusion), trop peu, et il
décroche (sous-stimulation, ennui)
\cite{sweller1988cognitive,sweller2019cognitive}. Les développements
ultérieurs distinguent charge intrinsèque (difficulté de la matière),
charge extrinsèque (présentation maladroite) et charge germane (effort
utile). Mayer, pour le multimédia, montre que le rythme, le texte et
l'image pèsent sur cette charge utile \cite{mayer2021multimedia}. Un
parcours unique, calé sur un « étudiant moyen », rate ces bascules :
c'est précisément ce qu'une adaptation par état cherche à éviter.

\subsection{États d'engagement : ennui, confusion, frustration, flow}

Quatre états, issus de \textsc{DAiSEE}, structurent la suite du mémoire
\cite{gupta2016daisee} :
\begin{itemize}
  \item L'engagement : implication positive, traitement en cours.
    on maintient ou on enrichit.
  \item L'ennui : sous-stimulation, amorce de
    retrait, on relance.
  \item La confusion : incompréhension encore active, on clarifie.
  \item La frustration : blocage, on soutient et on simplifie.
\end{itemize}
Chacun est annoté sur une intensité de 0 à 3. Distinguer ces états n'est
pas un luxe de taxonomie : chacun appelle une réponse différente.

Le flow de Csikszentmihalyi désigne le pôle productif : un
engagement soutenu quand le défi et les compétences s'équilibrent
\cite{csikszentmihalyi1990flow}. \textsc{DAiSEE} ne l'annote pas à part, 
l'état Engagement en est le proxy opératoire. On retient donc le flow
comme cadre théorique, et les quatre labels \textsc{DAiSEE} comme
taxonomie de détection.

D'Mello et Graesser suivent une séquence souvent observée en EIAH
(environnements informatiques pour l'apprentissage humain) : la
confusion, si elle se résout, aide (et peut ramener vers le flow), si elle
s'enlise, elle glisse vers la frustration, puis vers l'ennui
\cite{dmello2012dynamics}. Baker \emph{et al.} en concluent que
la frustration reste préférable à l'ennui, car la première suppose
un effort, le second un retrait \cite{baker2010better}. Cette réserve de
contexte compte : l'observation a été formulée dans des tuteurs
intelligents, plutôt que lors d'un visionnage isolé de cours en ligne.
Elle indique qu'un désengagement sans expression est, pédagogiquement,
plus grave qu'un blocage encore actif.

%==============================================================================
\section{IA générative pour l'éducation}
\label{sec:ea-genai}

Les tuteurs intelligents adaptent souvent un
contenu déjà écrit. Les grands modèles de langue rendent possible de le produire à la demande. Quatre axes sont examinés : génération
pédagogique, ingénierie de prompts, personnalisation, puis évaluation de
la qualité.

\subsection{LLMs et génération de contenu pédagogique}

Depuis GPT-3 et les modèles de la famille Llama, un grand modèle de langue
peut rédiger une explication, une analogie, un exercice ou un encouragement
en quelques secondes
\cite{brown2020language,touvron2023llama}. \textcite{kasneci2023chatgpt} ont très tôt
listé ce que cela change pour l'école : personnalisation possible, mais aussi
erreurs factuelles, biais, et décisions difficiles à retracer. Les revues qui ont suivi (Lo, Yan, Bond, entre
autres) ne contredisent pas ce constat
\cite{lo2023chatgpt,yan2024practical,bond2024systematic}. Laissé libre, le
modèle écrit « comme un cours » sans garantir que le cours soit juste.

Un garde-fou fréquent, distinct du prompt, est le RAG : on récupère d'abord
des extraits d'une base documentaire, puis on les fournit au modèle
\cite{lewis2020retrieval}. C'est ce mécanisme qu'\emolearn{} retient pour
coller le texte généré à des fiches de cours, plutôt que de le laisser
inventer une définition. Le LLM n'est donc pas, ici, un enseignant
autonome : c'est un rédacteur conditionné par une décision déjà prise
et par des documents vérifiés.

\subsection{Prompt engineering pour applications éducatives}

La qualité du texte dépend fortement du guidage : rôle, format, longueur,
exemples \cite{liu2023pretrain,white2023prompt,wei2022chain}. En éducation,
un prompt trop vague produit un cours générique, un prompt trop serré
peut brider l'utilité. Les patrons structurés permettent de conditionner
la génération selon l'état visé et le type d'intervention (clarifier,
exemplifier, exercer, soutenir).

Liu \emph{et al.} (2025) poussent le prompt dans une autre direction : non
pour écrire un cours, mais pour étiqueter un message (comportement,
émotion, cognition) \cite{liu2025agka}. Leur méthode AGKA ajoute au prompt
les définitions du guide d'annotation et quelques exemples. GPT-4 et
Llama~3 70B s'en sortent alors aussi bien que BERT sur des tâches binaires, 
dès que les classes se ressemblent (émotions épistémiques, phases CoI),
l'écart se creuse. On en retient surtout ceci : un LLM guidé peut aider à
classer du texte, il ne remplace pas un détecteur visage--voix.
\emolearn{} fait l'inverse : le modèle n'étiquette pas l'état, il rédige
le contenu après que la couche de perception a décidé, via une
bibliothèque de vingt patrons.

\subsection{Personnalisation et adaptation automatique}

Hu et Shao proposent PAGE (\emph{Psychologically-Aware Generative Education}),
un système GenAI qui adapte le style, la difficulté et les mécanismes de
soutien au profil Big Five et au feedback émotionnel en temps réel
\cite{hu2025page}. Sur 200 étudiants, PAGE améliore la satisfaction
émotionnelle (4,4/5 contre 3,6/5), le taux de complétion des tâches
(87,6\,\% contre 72,3\,\%) et la durée d'étude (+34\,\%) par rapport à un
système adaptatif classique, le bénéfice est particulièrement marqué pour
les étudiants à neuroticisme élevé (baisse de 48\,\% du décrochage). PAGE
est le travail le plus proche d'\emolearn{} du point de vue de la génération : il relie un signal affectif à un contenu produit par
LLM. L'écart se joue sur le quand et le comment. PAGE s'appuie sur
un profil Big Five (assez stable) et sur un retour émotionnel souvent
déclaré. Le présent travail vise l'état de ce clip : confusion
maintenant, ennui après trois minutes. PAGE ne décrit pas non plus un
moteur de règles nommé, ni un RAG sur fiches de cours : le style change,
mais on ne peut pas retracer « tel état $\rightarrow$ tel patron ».

Kong \emph{et al.} évaluent un autre mode d'adaptation : un outil de
représentation des discussions assisté par IA (représentations
comportementales, cognitives, ou mixtes) auprès de 122 étudiants
\cite{kong2025exploratory}. L'outil n'améliore pas significativement la
motivation, mais augmente la présence cognitive de haut niveau et la
performance. Il s'agit d'une médiation sur le forum, non d'une
génération de contenu de cours conditionnée par l'état d'un apprenant
isolé. Le résultat confirme qu'une personnalisation IA peut élever le
niveau cognitif, à condition d'en changer le canal (vidéo de cours plutôt
que fil de discussion).

\subsection{Évaluation de la qualité pédagogique}

La qualité du contenu généré se juge sur des critères humains (pertinence,
clarté, adéquation à l'état) complétés par des métriques automatiques de
forme (longueur, lisibilité, respect du patron). Aucune métrique automatique
ne se substitue à une validation experte ou utilisateur
\cite{kasneci2023chatgpt,hu2025page}. Les travaux de détection rapportent
surtout des F1 et des $\kappa$, les travaux génératifs rapportent des scores
de satisfaction. Un protocole multi-niveaux (perception, décision,
génération) reste rare. C'est l'un des choix méthodologiques d'\emolearn{}.

%==============================================================================
\section{Systèmes adaptatifs d'apprentissage}
\label{sec:ea-its}

Identifier un état ne suffit pas à en déduire une intervention. Les
tuteurs intelligents en sont le précédent le plus net : un modèle de
l'apprenant, une stratégie, un contenu. Trois étages sont distingués ici :
ce que sont les ITS, quelles stratégies pédagogiques ils mobilisent, puis
comment la décision est formalisée (règles et modèles).

\subsection{Intelligent Tutoring Systems (ITS)}

Les systèmes de tutorat intelligent (ITS) combinent historiquement un modèle
de l'apprenant, un modèle pédagogique et une interface adaptative
\cite{vanlehn2011relative}. VanLehn montre que les meilleurs ITS approchent
l'efficacité d'un tuteur humain, à condition que le domaine soit bien
modélisé \cite{vanlehn2011relative}. Le contenu est le plus souvent
pré-écrit : chaque rétroaction a été rédigée à l'avance pour un état
de connaissances prévu.

Deux familles d'exemples illustrent ce schéma sans en changer la
logique. Les tuteurs cognitifs, dont la revue de VanLehn rend compte,
adaptent surtout la prochaine étape d'un parcours : l'item suivant,
un indice, une remédiation calée sur ce que l'apprenant maîtrise encore
mal. Les tuteurs affectifs, dont AutoTutor et les dynamiques décrites par
D'Mello et Graesser sont les références les plus citées, observent aussi
la confusion ou la frustration pendant le dialogue, puis ajustent
l'étayage \cite{dmello2012dynamics,graesser2020emotions}. Dans les deux cas,
l'intervention reste choisie dans un répertoire d'auteurs, pas rédigée à
la demande.

Leur limite, dans un MOOC aux contenus ouverts et aux états affectifs
fluctuants, est double. D'une part, on ne peut pas pré-écrire un texte pour
chaque combinaison d'état et de concept. D'autre part, le modèle de
l'apprenant est surtout cognitif (ce qu'il sait), rarement émotionnel (s'il
s'ennuie ou s'il est bloqué), même si des travaux récents visent une
explication personnalisée des décisions de tuteurs \cite{conati2021toward}.
C'est cette limite que les modèles de langue visent à lever : générer le
contenu, tout en gardant une décision lisible.

\subsection{Stratégies d'adaptation pédagogique}

Les stratégies d'adaptation classiques ne portent pas sur quel
algorithme détecte l'état, mais sur l'action à entreprendre une fois
l'état connu. On distingue notamment la simplification (réduire
la charge, Sweller), l'enrichissement (approfondir lorsque
l'apprenant est déjà engagé), la remédiation (revenir sur un
point mal compris) et l'encouragement (soutenir un effort qui
s'enlise) \cite{sweller1988cognitive,vanlehn2011relative}.

Ces familles se recoupent avec les quatre états \textsc{DAiSEE} : la
confusion appelle plutôt une clarification, la frustration un soutien et
une simplification, l'ennui une relance, l'engagement un maintien ou un
enrichissement. PAGE adapte plutôt le style et la
difficulté à un profil de personnalité \cite{hu2025page}, les ITS
adaptent surtout la prochaine étape d'un parcours cognitif. Dans
les deux cas, une stratégie sans contenu à produire (ou sans règle pour
la choisir) reste un principe, pas un système.

\subsection{Règles et modèles de décision}

Deux grandes manières de décider coexistent. Les modèles
cognitifs (traçage bayésien des connaissances, théorie des espaces de
connaissances) estiment ce que l'apprenant maîtrise, puis sélectionnent
un item \cite{vanlehn2011relative}. Les règles explicites
relient un état observé à une action nommée, au prix d'une rigidité, mais
avec une traçabilité que Conati \emph{et al.} jugent utile pour expliquer
un tuteur \cite{conati2021toward}.

\emolearn{} se place délibérément du côté des règles. Une matrice
\emph{état\,$\rightarrow$\,type de contenu} (explication, exemple,
exercice, encouragement) fixe le quoi, des règles temporelles sur
les séquences de clips fixent le quand (un épisode isolé n'a pas le
même poids qu'une frustration qui dure). Le moteur reste interprétable :
chaque adaptation doit être justifiable par l'état détecté, la stratégie
retenue et le patron mobilisé. Le détail de cette matrice appartient au
chapitre~\ref{chap:conception}, l'état de l'art retient surtout le
principe : sans règle nommée, la génération n'est plus une décision
pédagogique, c'est une improvisation.

%==============================================================================
\section{Synthèse critique et positionnement}
\label{sec:ea-synthese}

Cette section clôt l'état de l'art par une lecture critique des
travaux (ce qu'ils établissent, ce qu'ils laissent ouvert), puis par le
positionnement d'\emolearn{}. Le tableau~\ref{tab:ea-comparatif} en
est l'instrument : référence, approche / modèle, contexte, forces, limites,
résultats clés. Ce tableau compare des méthodes. Il se lit
par blocs (cadres et corpus, puis détection, ITS et génération), plutôt
que comme un classement unique.

\providecommand{\eatabhead}{}
\renewcommand{\eatabhead}{%
\rowcolor{gray!20}
\textbf{Référence} & \textbf{Approche / modèle} & \textbf{Contexte} &
\textbf{Forces} & Limites & Résultats clés \\}

Le tableau~\ref{tab:ea-comparatif} présente la synthèse critique des travaux retenus.

{\scriptsize
\setlength{\tabcolsep}{2.4pt}
\setlength{\LTcapwidth}{\textwidth}
\renewcommand{\arraystretch}{1.08}
\begin{xltabular}{\textwidth}{|>{\raggedright\arraybackslash}X|>{\raggedright\arraybackslash}X|>{\raggedright\arraybackslash}X|>{\raggedright\arraybackslash}X|>{\raggedright\arraybackslash}X|>{\raggedright\arraybackslash}X|}
\hline
\eatabhead
\hline
\endfirsthead
\hline
\eatabhead
\hline
\endhead
\endfoot
\endlastfoot
\rowcolor{gray!12}
\multicolumn{6}{|l|}{\textbf{Cadres théoriques, enjeux MOOC et corpus}} \\
\hline
\textcite{kizilcec2013deconstructing}, \textcite{hew2016promoting}
  & Clustering de trajectoires, études de cas
  & Engagement dans les MOOC
  & Identification de sous-populations et de leviers
  & Proxies comportementaux (clics, progression)
  & L'engagement explique une part majeure du décrochage \\
\hline
\textcite{fredricks2004school}
  & Synthèse conceptuelle
  & Psychologie de l'éducation
  & Trois dimensions (comportement, cognition, émotion)
  & Pas d'estimateur automatique
  & Cadre de référence encore dominant \\
\hline
\textcite{garrison1999critical}
  & Community of Inquiry (CoI)
  & Discussions en ligne
  & Quatre phases de présence cognitive
  & Conçu pour les forums, pas pour la vidéo
  & Schéma d'annotation largement repris \\
\hline
\textcite{chi2014icap}
  & Taxonomie ICAP
  & Engagement cognitif
  & Lien mode d'activité $\rightarrow$ apprentissage
  & Codage d'activités, pas d'affect facial/vocal
  & Plus l'activité est constructive, plus le traitement est profond \\
\hline
\textcite{picard1997affective}
  & Informatique affective
  & Reconnaissance d'émotions
  & Fondement théorique de la détection
  & Socle, sans pipeline pédagogique
  & L'affect est mesurable et actionnable \\
\hline
\textcite{pekrun2006control}
  & Théorie contrôle--valeur
  & Émotions d'accomplissement
  & Lien valence / activation / performance
  & Pas d'opérationnalisation logicielle
  & Les émotions modulent attention et persévérance \\
\hline
\textcite{dmello2012dynamics}, \textcite{baker2010better}
  & Dynamique d'états affectifs
  & Apprentissage complexe médiatisé
  & Séquence confusion--frustration--ennui
  & Pas de génération de contenu
  & Mieux vaut être frustré qu'ennuyé \\
\hline
\textcite{dmello2015review}
  & Méta-analyse
  & États affectifs en EIAH
  & Incidence relative des états, fusion multimodale
  & Synthèse, pas de système
  & Ennui, confusion, engagement et frustration sont les états clés \\
\hline
\textcite{sweller1988cognitive}
  & Théorie de la charge cognitive
  & Conception pédagogique
  & Explique surcharge et sous-stimulation
  & Théorie, sans détecteur
  & Ajuster la complexité au profil de l'apprenant \\
\hline
\textsc{DAiSEE} \cite{gupta2016daisee}
  & Corpus vidéo annoté
  & Engagement \og dans la nature \fg{}
  & Quatre états, intensité 0--3
  & Fort déséquilibre de classes
  & Taxonomie opératoire pour la détection faciale \\
\hline
\textsc{EmotiW} \cite{dhall2018emotiw}
  & Défi d'affect \og dans la nature \fg{}
  & Engagement étudiant, vidéo
  & Conditions peu contrôlées, piste engagement
  & Niveaux d'engagement, pas les 4 états \textsc{DAiSEE}
  & Référence de challenge, non retenue comme taxonomie \\
\hline
\textsc{IEMOCAP} \cite{busso2008iemocap}
  & Corpus audiovisuel
  & Émotion conversationnelle
  & Voix riche, visage et voix alignés
  & Taxonomie émotionnelle, pas d'engagement MOOC
  & Référence SER, cadre de fusion alignée \\
\hline
\rowcolor{gray!12}
\multicolumn{6}{|l|}{\textbf{Détection de l'engagement, ITS et génération adaptative}} \\
\hline
\textcite{sakeef2025detecting}
  & Revue CoI / ICAP / KCSA + TALN
  & Forums de cours en ligne
  & Cartographie des cadres et des méthodes IA
  & Texte seul, pas de webcam/micro
  & L'automatisation du codage est indispensable au suivi continu \\
\hline
\textcite{hu2021automatic}
  & CNN + XAI
  & Présence cognitive (CoI)
  & Explicabilité au niveau des mots
  & Texte seul, généralisabilité limitée
  & $\kappa$ de Cohen $= 0{,}528$, comparable aux Random Forests \\
\hline
\textcite{liu2021exploring}
  & SNA + ENA
  & Forum MOOC (3\,925 apprenants)
  & Lien réseau social / traitement cognitif / réussite
  & Analyse collective, hors temps réel
  & Centralité et connexions cognitives plus fortes chez les performants \\
\hline
\textcite{liu2022automated}
  & BERT-CNN
  & Forums MOOC (8\,867 apprenants)
  & Détection conjointe émotion + cognition
  & Pas d'intervention générée
  & F1 $+10\,\%$ (émotion) et $+8\,\%$ (cognition), effet combiné sur la réussite \\
\hline
\textcite{liu2023moocbert}
  & MOOC-BERT (pré-entraînement multi-cours)
  & Présence cognitive dans les MOOC
  & Généralisation inter-cours, sans ingénierie de traits
  & Diagnostic enseignant, sans contenu adaptatif
  & Surpasse ML et DL représentatifs, analyse longitudinale des phases \\
\hline
\textcite{liu2025agka}
  & LLM + AGKA (guides d'annotation)
  & Classification d'engagement textuel
  & GPT-4 / Llama~3 rivalisent avec BERT en binaire
  & LLM comme détecteur, pas comme générateur ancré
  & Retrait sur les taxonomies multi-classes sémantiquement proches \\
\hline
\textcite{whitehill2014faces}
  & Classifieur d'expressions faciales
  & Engagement en classe / e-learning
  & Fiabilité comparable à l'annotation humaine
  & Pas de voix, pas de génération
  & L'engagement facial est mesurable automatiquement \\
\hline
\textcite{xiong2024cnntransformer}
  & CNN-Transformer (visage + posture)
  & Classe réelle annotée
  & Combine grain fin (visage) et grain large (posture)
  & Pas de voix, pas de boucle pédagogique
  & Accuracy $92{,}9\,\%$, supérieure aux baselines DL \\
\hline
\textcite{baevski2020wav2vec2}
  & wav2vec~2.0
  & Représentations de parole
  & Encodeur auto-supervisé réutilisable (SER)
  & Générique, pas d'engagement pédagogique
  & Socle de fine-tuning vocal largement adopté \\
\hline
\textcite{chen2022wavlm}
  & WavLM
  & Parole auto-supervisée
  & Représentations plus riches que wav2vec~2.0 (benchmark SUPERB)
  & Pas d'engagement pédagogique natif
  & Encodeur retenu par \emolearn{} pour le SER \\
\hline
\textcite{vanlehn2011relative}
  & ITS, contenu pré-écrit
  & Tutorat cognitif
  & Efficacité proche d'un tuteur humain
  & Peu d'affect, contenu figé
  & Les meilleurs ITS approchent le tutorat humain \\
\hline
\textcite{brown2020language}, \textcite{touvron2023llama}
  & LLM (GPT-3, Llama)
  & Génération de texte
  & Production à la demande, few-shot
  & Hallucinations, pas d'état d'engagement
  & Faisabilité d'un contenu pédagogique généré \\
\hline
\textcite{kasneci2023chatgpt}
  & Revue LLM pour l'éducation
  & Enseignement et apprentissage
  & Cartographie opportunités / risques
  & Pas de prototype multimodal
  & Fiabilité, biais et traçabilité restent ouverts \\
\hline
\textcite{lewis2020retrieval}
  & RAG (index + générateur)
  & Tâches NLP intensives en connaissances
  & Ancrage documentaire, moins d'hallucinations
  & Non couplé à l'engagement
  & Le contexte récupéré conditionne la génération \\
\hline
\textcite{kong2025exploratory}
  & Outil de représentation IA
  & Apprentissage collaboratif en ligne ($n=122$)
  & Hausse de la présence cognitive de haut niveau
  & Canal forum, pas de contenu de cours
  & Performance améliorée, motivation non significativement accrue \\
\hline
\textcite{hu2025page}
  & GenAI PAGE (Big Five + émotion)
  & Contenu pédagogique personnalisé ($n=200$)
  & Adaptation du style, de la difficulté et du soutien
  & Pas de fusion visage--voix, ni RAG, ni règles explicites
  & Satisfaction $4{,}4/5$ vs $3{,}6$, complétion $87{,}6\,\%$ vs $72{,}3\,\%$ \\
\hline
\end{xltabular}
\xlcaption{Synthèse critique de l'état de l'art.}
\label{tab:ea-comparatif}
}

Aucun des travaux du tableau ne ferme à lui seul la boucle qui intéresse
ce mémoire.

\subsection*{Lecture critique : cinq lacunes}

Cinq limites se recoupent.

\begin{enumerate}[leftmargin=1.4em]
  \item \textbf{Signal:} Les détecteurs récents lisent surtout des forums
  \cite{liu2022automated,liu2023moocbert}. Ils laissent de côté
  l'apprenant qui visionne sans poster.
  \item \textbf{Décision:} Quand une intervention apparaît, elle règle le
  style d'un LLM (PAGE) ou rend visible un forum (Kong), rarement une
  règle nommée « état $\rightarrow$ type de contenu ».
  \item \textbf{Génération:} Le RAG ancre des faits
  \cite{lewis2020retrieval} sans connaître l'ennui ou la confusion du
  moment, la génération conditionnée par un état facial et vocal reste
  rare.
  \item \textbf{Évaluation:} On mesure soit un F1, soit une satisfaction :
  presque jamais perception, décision et texte généré sur le même banc.
  \item \textbf{Modalités:} Les approches visuelles
  \cite{whitehill2014faces,xiong2024cnntransformer} estiment
  l'engagement sur le visage, sans intégrer, le plus souvent, le
  signal vocal.
\end{enumerate}

\subsection{Positionnement d'EmoLearn}
\label{subsec:ea-positionnement}

Le positionnement d'\emolearn{} répond à ces cinq lacunes. L'objet du
mémoire est une chaîne unique : percevoir un état à partir du visage et
de la voix, décider explicitement s'il faut intervenir et avec quel type
de contenu, puis produire un texte ancré documentairement
(chapitre~\ref{chap:conception}). Cette chaîne vise l'apprenant qui
visionne une vidéo sans contribution écrite, le signal exploité est
donc facial et vocal.

Par rapport aux détecteurs textuels (BERT-CNN, MOOC-BERT, AGKA), le
changement porte sur le canal, et non sur une amélioration de F1. Par
rapport aux classifieurs visuels
\cite{whitehill2014faces,xiong2024cnntransformer}, trois maillons
supplémentaires sont introduits : une voie vocale, une fusion dont le
désaccord reste lisible, et une décision qui déclenche une génération.
Les ITS de \textcite{vanlehn2011relative} montrent qu'une médiation peut
approcher le tutorat humain, avec toutefois un contenu figé et une prise
en compte limitée de l'affect. Dans \emolearn{}, l'état détecté
conditionne un contenu produit à la demande. PAGE \cite{hu2025page}
adapte style et difficulté à partir d'un profil et d'une émotion, sans
fusion visage--voix ni règle nommée. \textcite{kong2025exploratory}
améliore la présence cognitive sur un forum, et non le contenu de cours.
Le mémoire se situe entre ces deux familles, avec vingt patrons A1--D5,
une cascade inspectable «~état~$\rightarrow$~type~», un RAG puis un LLM.

Ces choix s'accompagnent de limites, énoncées avant la conception. La
chaîne est éprouvée en simulation sur corpus, et non déployée dans un
MOOC réel. En production, le visage \textsc{DAiSEE} et la voix
\textsc{IEMOCAP} ne proviennent pas du même locuteur, seul un
sous-ensemble aligné sert à mesurer la fusion. Le mapping
émotion~$\rightarrow$~engagement constitue une hypothèse de construit,
distincte d'une équivalence psychologique. L'évaluation du
chapitre~\ref{chap:validation} porte sur la perception, la décision et le
texte généré, sans étude longitudinale d'apprenants.

La contribution consiste à relier l'informatique affective, les tuteurs
intelligents et les LLM ancrés sur le cas d'un visionnage individuel, de
manière assez explicite pour reconstruire chaque étape.

%==============================================================================
\section{Conclusion}
\label{sec:ea-concl}

La littérature documente un décrochage massif des MOOC, dans lequel
l'engagement occupe une place centrale. Les cadres CoI et ICAP éclairent
l'analyse des forums, ils s'appliquent peu à un visionnage sans
message. Pekrun, D'Mello et Sweller relient émotion et charge cognitive
à la suite du parcours, Baker précise que cette observation provient
surtout de tuteurs intelligents. Les détecteurs de texte (BERT-CNN,
MOOC-BERT, AGKA) et les détecteurs de visage (Whitehill, Xiong,
\textsc{DAiSEE}) atteignent des performances utiles, sous réserve des
biais de corpus. ITS, PAGE et l'outil de Kong montrent qu'une médiation
par IA peut modifier l'expérience d'apprentissage, rarement toutefois à
partir d'un visage et d'une voix, avec une décision relisible.

Ces travaux laissent ouverte la conception d'une chaîne dont chaque
étape peut être reconstruite. Le chapitre suivant en propose la
formalisation.
